\documentclass[11pt]{article}

\usepackage[utf8]{inputenc}
\usepackage[T1]{fontenc}
\usepackage{lmodern}
\usepackage{microtype}
\usepackage[margin=1in]{geometry}
\usepackage{amsmath,amssymb,amsthm}
\usepackage{graphicx}
\usepackage{booktabs}
\usepackage{xcolor}
\usepackage[numbers,sort&compress]{natbib}
\usepackage[colorlinks=true,linkcolor=blue!50!black,citecolor=green!40!black,urlcolor=blue!60!black]{hyperref}
\usepackage[capitalize]{cleveref}
\usepackage{enumitem}

\newtheorem{proposition}{Proposition}
\newcommand{\todo}[1]{{\color{red!75!black}\textbf{[TODO:} #1\textbf{]}}}
\newcommand{\R}{\mathbb{R}}
\newcommand{\E}{\mathbb{E}}
\newcommand{\cS}{\mathcal{S}}
\newcommand{\cA}{\mathcal{A}}
\newcommand{\cO}{\mathcal{O}}
% Include a generated table/figure if it exists, otherwise show a placeholder.
\newcommand{\maybeinput}[1]{\IfFileExists{#1}{\input{#1}}{\fbox{\parbox{0.8\linewidth}{\centering\small\todo{run \texttt{scripts/aggregate.py} to generate \texttt{#1}}}}}}
\newcommand{\maybegraphic}[2]{\IfFileExists{#2}{\includegraphics[#1]{#2}}{\fbox{\parbox[c][4cm]{0.8\linewidth}{\centering\small\todo{run \texttt{scripts/aggregate.py} to generate \texttt{#2}}}}}}

\title{Memory, Layout Diversity and Reward Shaping in\\Partially Observable Maze Foraging:\\A Reproducible Small-Scale Study}
\author{Bishwaswarup\thanks{OUSUMS Lab, Department of Physics, Indian Institute of Science, Bangalore, India. \texttt{bishwaswarup@iisc.ac.in}}
% \and Co-author \thanks{...}
}
\date{Preprint --- \today}

\begin{document}
\maketitle

% =====================================================================
\begin{abstract}
We study a deliberately small foraging task. An agent with a $7\times7$
egocentric view must reach a food source in procedurally generated
$11\times11$ perfect mazes. The task is cheap enough to train many agents
on a laptop CPU, yet it exposes three design choices that matter widely in
deep reinforcement learning: \emph{memory} (none, frame stacking, or an
LSTM), \emph{training-layout diversity} (from a single maze to a fresh maze
every episode), and \emph{reward design} (sparse, exact potential-based
shaping, or an ad-hoc shaped reward). Every agent is trained with PPO and
evaluated on the same set of 500 held-out layouts, and on larger mazes, with
success rate, SPL and bootstrap confidence intervals over seeds. As
references we include a random walk, an oracle, and a hand-coded
right-hand wall follower. The wall follower uses one integer of memory and
always succeeds in a perfect maze.
\todo{One or two sentences stating the main empirical findings once the experiments are run, e.g.\ how many layouts are needed before held-out success matches training success, whether recurrence beats frame stacking, and whether learned agents beat the wall follower on SPL.}
We do not propose a new method. Our aim is a transparent, fully
reproducible reference point, and we release the code, configurations and
evaluation protocol.
\end{abstract}

% =====================================================================
\section{Introduction}

Social insects forage efficiently with tiny brains and limited
perception~\citep{wehner2003desert,bonabeau1999swarm}. This makes ant-like
foraging a natural testbed for learning agents that must act on local,
partial information. In deep reinforcement learning (RL), even the simplest
version of the task, one agent looking for food in a maze, touches three
recurring questions:

\begin{enumerate}[leftmargin=*,itemsep=2pt]
  \item \textbf{How much memory does a partially observed agent need?} Frame
  stacking~\citep{mnih2015human} and recurrent policies~\citep{hausknecht2015drqn,ni2022recurrent}
  are the standard choices, and they trade simplicity for capacity differently.
  \item \textbf{How many distinct environments must an agent see to
  generalise?} Agents trained on few procedurally generated levels tend to
  memorise them~\citep{zhang2018overfitting,cobbe2019quantifying,cobbe2020procgen,kirk2023survey}.
  \item \textbf{Does reward shaping help, and is it safe?} Potential-based
  shaping provably preserves optimal policies~\citep{ng1999policy}. Ad-hoc
  shaping terms do not, and they are common in practice.
\end{enumerate}

We study all three on one small, fully specified task. Our contributions
are empirical and practical:
\begin{itemize}[leftmargin=*,itemsep=2pt]
  \item an open, seeded Gymnasium~\citep{towers2024gymnasium} environment with
  configurable layout pools and reward presets, plus a test suite that checks
  maze validity, seeding and the telescoping property of the shaping term;
  \item a fixed evaluation protocol: held-out layouts, success and SPL~\citep{anderson2018evaluation},
  bootstrap CIs over seeds~\citep{agarwal2021precipice,henderson2018matters},
  no checkpoint selection on evaluation data, and simple non-learned baselines;
  \item controlled experiments on memory $\times$ layout diversity, reward
  design, and field of view, with zero-shot transfer to larger mazes.
\end{itemize}

% =====================================================================
\section{Related Work}

\paragraph{Partial observability.} Maze navigation from local views is a
POMDP~\citep{kaelbling1998pomdp}. Frame stacking gives a fixed-length
history~\citep{mnih2015human}, while recurrent policies~\citep{hochreiter1997lstm,hausknecht2015drqn}
can in principle integrate unbounded history.
\citet{ni2022recurrent} show that well-tuned recurrent model-free RL is a
strong POMDP baseline, and POPGym~\citep{morad2023popgym} benchmarks memory
models across many tasks. Navigation agents in richer 3D mazes also rely on
recurrence~\citep{mirowski2017navigate}.

\paragraph{Generalisation and procedural generation.}
\citet{zhang2018overfitting} and \citet{cobbe2019quantifying} show that
deep RL agents overfit to small sets of training levels, and Procgen~\citep{cobbe2020procgen}
made the number of training levels a standard axis of evaluation; see~\citet{kirk2023survey}
for a survey. Grid-world suites such as Minigrid~\citep{chevalier2023minigrid}
include maze tasks similar to ours. Our environment is intentionally simpler,
so that many seeds and configurations fit on a CPU.

\paragraph{Reward shaping.} Potential-based shaping $F=\gamma\Phi(s')-\Phi(s)$
leaves the optimal policy unchanged~\citep{ng1999policy} and is equivalent to
value initialisation~\citep{wiewiora2003potential}. In episodic tasks the
potential of terminal states must be handled with care~\citep{grzes2017reward}.

\paragraph{Ant-inspired algorithms.} Ant colony optimisation~\citep{dorigo2004aco}
and swarm intelligence~\citep{bonabeau1999swarm} use stigmergy, coordination
through the environment. Our single-agent study is the first step towards a
multi-agent setting with learned pheromone use (\cref{sec:discussion}).

% =====================================================================
\section{Task and Environment}\label{sec:env}

\begin{figure}[t]
  \centering
  \maybegraphic{width=0.85\linewidth}{figures/env.pdf}
  \caption{\textbf{The task.} Left: an $11\times11$ perfect maze with the nest
  (brown), food (green), the ant (red) and its $7\times7$ field of view (blue
  tint). Right: the three observation channels the agent receives.}
  \label{fig:env}
\end{figure}

\paragraph{Mazes.} A layout is an $11\times11$ grid whose border is wall. It
is a \emph{perfect} maze: a spanning tree over the $5\times5$ lattice of
odd-coordinate ``rooms'', generated by randomised depth-first search (which
favours long corridors and is not uniform over spanning trees). Exactly one
simple path joins any two of its $49$ free cells. A layout is a
deterministic function of an integer seed. Training layouts come either
from a fixed pool of $N$ seeds ($N\in\{1,8,64,512\}$) or are freshly sampled
every episode ($N=\infty$). The $5\times5$ room lattice has
$557{,}568{,}000$ spanning trees, but randomised DFS samples them far from
uniformly: of $200{,}000$ procedurally sampled layouts only $158{,}430$ were
distinct, and $255$ samples coincided with one of $500$ layouts drawn from a
disjoint \emph{seed} range (\texttt{scripts/layout\_stats.py}). Separating
seeds is therefore not enough. We assign every layout to a split by a hash of
its cells, with one layout in ten reserved for testing. Training only ever uses
train-split layouts and held-out evaluation only test-split layouts, so the
overlap is zero by construction (verified: $0$ hits in the same $200{,}000$ samples).

\paragraph{Episodes.} The nest and food are placed uniformly at random
among free cells, subject to a geodesic (shortest-path) distance of at least
3. Placements are re-sampled every episode even when the layout pool is
fixed. The agent starts at the nest. An episode terminates when the agent
reaches the food and is truncated after $T_{\max}=6\times49=294$ steps.

\paragraph{POMDP.} The state $s=(\text{layout},p_{\text{agent}},p_{\text{nest}},p_{\text{food}})$
is hidden. Actions are $\cA=\{\uparrow,\rightarrow,\downarrow,\leftarrow\}$,
and a move into a wall leaves the agent in place. The observation $o\in\{0,1\}^{3\times7\times7}$
is an egocentric window of radius $k=3$ with three channels: wall (cells
outside the grid count as walls), food, and nest (\cref{fig:env}). The window
covers about 40\% of the grid but rarely shows the food at the start of an
episode, so the agent has to explore.

\paragraph{Rewards.} We compare three presets.
\emph{Sparse}: $r=\mathbb{1}[\text{food reached}]$.
\emph{PBRS}: sparse plus $F(s,s')=\gamma\Phi(s')-\Phi(s)$ with
\begin{equation}
  \Phi(s) = -\,\frac{d(p_{\text{agent}},p_{\text{food}})}{D},\qquad D=\max_{c} d(c,p_{\text{food}}),
\end{equation}
where $d$ is the geodesic (BFS) distance in the maze. We set $\Phi\equiv0$ at
the goal.
\emph{Legacy}: the reward used in an earlier version of this code:
$-0.01$ per step, $-0.1$ per wall bump, $0.15\,\Delta d_{1}$ where $d_1$ is
the \emph{Manhattan} distance (with no $\gamma$ factor), a revisit penalty
$-0.05\min(n_{\text{visits}}-1,5)$, and $+10$ on success.

\begin{proposition}\label{prop:pbrs}
For any trajectory $s_0,\dots,s_\tau$ that ends at the goal or by truncation,
the discounted PBRS return equals the discounted sparse return plus
$\gamma^{\tau}\Phi(s_\tau)-\Phi(s_0)$. $\Phi(s_0)$ is fixed by the
episode's start state and $\Phi(s_\tau)=0$ at the goal, so
on episodes that reach the goal the two rewards rank policies identically~\citep{ng1999policy,grzes2017reward}.
\end{proposition}
\begin{proof}
$\sum_{t=0}^{\tau-1}\gamma^{t}\big(\gamma\Phi(s_{t+1})-\Phi(s_t)\big)$ telescopes
to $\gamma^{\tau}\Phi(s_\tau)-\Phi(s_0)$.
\end{proof}

The identity is checked numerically in the test suite. The legacy reward is
not of this form: the revisit penalty depends on history, and the Manhattan
term ignores walls and $\gamma$. It can therefore change which policies are
optimal.

% =====================================================================
\section{Agents and Training}\label{sec:agents}

All agents share one feature extractor: two $3\times3$ convolutions (16 and 32
channels, ReLU), adaptive average pooling to $3\times3$, and a 64-unit linear
layer. Separate two-layer MLP heads (64 units) produce the policy and the
value. The memory variants differ only in their input or core:
\begin{itemize}[leftmargin=*,itemsep=1pt]
  \item \textbf{None}: the current observation only (a reactive policy);
  \item \textbf{Frame stack}: the last 4 observations concatenated along channels~\citep{mnih2015human};
  \item \textbf{LSTM}: a 64-unit LSTM~\citep{hochreiter1997lstm} after the extractor (RecurrentPPO from SB3-Contrib).
\end{itemize}
We train with PPO~\citep{schulman2017ppo} and GAE~\citep{schulman2016gae} in
Stable-Baselines3~\citep{raffin2021sb3}, using one hyperparameter setting
for every condition (\cref{tab:hparams}) and no per-condition tuning. Each run
lasts $10^6$ environment steps.
\todo{State whether a small pilot sweep was used to choose the shared hyperparameters; see \citet{andrychowicz2021matters}.}

\begin{table}[t]
\centering\small
\caption{Hyperparameters, shared by all conditions (\texttt{configs/default.yaml}).}
\label{tab:hparams}
\begin{tabular}{@{}ll@{\hspace{2em}}ll@{}}
\toprule
Parallel envs & 8 & Discount $\gamma$ & 0.99 \\
Rollout length / env & 256 & GAE $\lambda$ & 0.95 \\
Minibatch size & 256 & Clip $\epsilon$ & 0.2 \\
Epochs / update & 10 & Entropy coef. & 0.01 \\
Learning rate & $3\times10^{-4}$ (Adam) & Value coef. & 0.5 \\
Total env steps & $10^6$ & Max grad norm & 0.5 \\
Frame stack & 4 & LSTM hidden & 64 \\
\bottomrule
\end{tabular}
\end{table}

% =====================================================================
\section{Evaluation Protocol}\label{sec:protocol}

\begin{itemize}[leftmargin=*,itemsep=2pt]
  \item \textbf{Held-out set}: 500 episodes on 500 test-split layouts
  (\cref{sec:env}), which never occur in training, each with a fixed nest
  and food placement. Every method sees exactly the
  same episodes, so comparisons are paired.
  \item \textbf{Training distribution}: 200 episodes on training layouts
  with fresh placements. The difference between training and held-out success
  is the generalisation gap.
  \item \textbf{Size transfer}: 500 held-out episodes on $15\times15$ and $21\times21$
  mazes with the same $7\times7$ window, made possible by adaptive pooling.
  \item \textbf{Metrics}: success rate, and SPL $=\frac{1}{M}\sum_i S_i\,\ell_i/\max(p_i,\ell_i)$~\citep{anderson2018evaluation},
  where $\ell_i$ is the shortest-path length and $p_i$ the agent's path length.
  \item \textbf{Action selection}: we report both stochastic (sampling) and
  deterministic (arg-max) policies. A deterministic reactive policy can
  cycle forever in a maze, so the two can differ substantially.
  \item \textbf{Statistics}: 5 seeds per condition. We report means with 95\%
  percentile-bootstrap CIs over seeds~\citep{agarwal2021precipice}. We
  evaluate the \emph{final} policy; held-out data is never used for checkpoint selection.
  \item \textbf{Baselines}: a uniform random walk; a right-hand \emph{wall
  follower} that sees only the four adjacent cells and keeps one heading
  variable, and that provably visits every cell of a perfect maze and so
  always succeeds; and an \emph{oracle} that follows the shortest path.
\end{itemize}

\begin{table}[t]
\centering\small
\caption{Non-learned baselines on the held-out episodes (500 layouts per size).
The wall follower always succeeds, but its path is on average much longer than the shortest path.}
\label{tab:baselines}
\maybeinput{tables/baselines.tex}  % generated from results/baselines.json by scripts/aggregate.py
\end{table}

% =====================================================================
\section{Results}\label{sec:results}

\todo{All numbers below are generated by \texttt{scripts/aggregate.py}.
Write the narrative only after the runs finish. Report what the data shows,
including null or negative results.}

\subsection{Memory $\times$ layout diversity}

\begin{figure}[t]
  \centering
  \maybegraphic{width=\linewidth}{figures/main_generalisation.pdf}
  \caption{\textbf{Generalisation gap.} Success (stochastic policy) on
  training layouts (left) and held-out layouts (right) against the number of
  distinct training layouts. Shaded regions: 95\% bootstrap CI over 5 seeds.
  Dashed and dotted lines: wall follower and random walk.}
  \label{fig:main}
\end{figure}

\begin{table}[t]
\centering\scriptsize
\caption{Held-out performance (\%) by memory type and number of training
layouts ($\infty$ = fresh layout every episode). Mean [95\% CI] over 5 seeds.}
\label{tab:main}
\maybeinput{tables/main.tex}
\end{table}

\todo{Describe \cref{fig:main} and \cref{tab:main}. Suggested questions:
(i) at which pool size does held-out success approach training success?
(ii) does the ranking of memory types change with diversity?
(iii) do deterministic and stochastic evaluation differ more for the reactive (no-memory) agent?
(iv) how do learned agents compare with the wall follower on SPL?}

\begin{figure}[t]
  \centering
  \maybegraphic{width=0.55\linewidth}{figures/main_curves.pdf}
  \caption{Held-out success during training ($N=\infty$). Mean $\pm$ s.d. over seeds.}
  \label{fig:curves}
\end{figure}

\subsection{Reward design}

\begin{figure}[t]
  \centering
  \maybegraphic{width=\linewidth}{figures/reward.pdf}
  \caption{Held-out success and SPL for the three reward presets with 8 and
  $\infty$ training layouts (frame-stack agent). Legacy reward with 8 layouts
  reproduces the training setup of the earlier code release (\cref{app:errata}).}
  \label{fig:reward}
\end{figure}

\begin{table}[t]
\centering\scriptsize
\caption{Reward ablation, held-out (\%).}
\label{tab:reward}
\maybeinput{tables/reward.tex}
\end{table}

\todo{Describe \cref{fig:reward}: does PBRS speed learning or improve the
final policy compared with sparse? Does the legacy reward lead to
qualitatively different behaviour (e.g.\ fewer revisits but lower SPL)?}

\subsection{Field of view}

\begin{figure}[t]
  \centering
  \maybegraphic{width=\linewidth}{figures/window.pdf}
  \caption{Held-out performance against window radius $k\in\{1,2,3\}$ ($3\times3$ to $7\times7$ views), $N=\infty$.}
  \label{fig:window}
\end{figure}

\todo{Describe \cref{fig:window}. Does memory matter more as the view shrinks?}

\subsection{Transfer to larger mazes}

\begin{table}[t]
\centering\small
\caption{Held-out success (\%, stochastic) on larger mazes, zero-shot.}
\label{tab:transfer}
\maybeinput{tables/main_transfer.tex}
\end{table}

\todo{Describe \cref{tab:transfer}, comparing with the baselines in \cref{tab:baselines}.}

% =====================================================================
\section{Discussion and Limitations}\label{sec:discussion}

\todo{Summarise the findings in 1--2 paragraphs.}

\paragraph{Limitations.} The environment is small and the observation is
symbolic. Mazes are perfect (tree-structured), which is exactly the setting
where a wall follower is complete. Braided mazes with loops would remove
that guarantee. We use one hyperparameter setting for all conditions, so a
condition that would benefit from different settings (e.g.\ longer
sequences for the LSTM) may be disadvantaged. Five seeds give wide
intervals for some conditions. Compute was limited to CPUs.

\paragraph{Towards collective foraging.} The repository contains
prototype components for a multi-agent extension (a diffusing,
evaporating pheromone field and inter-agent message channels), which we
leave to future work.

% =====================================================================
\section*{Reproducibility Statement}
Code, configurations and the evaluation protocol are available at
\url{https://github.com/Bishwaswarup/ANTS_RL}. Each figure and table in
\cref{sec:results} is produced by \texttt{scripts/run\_experiments.py} followed
by \texttt{scripts/aggregate.py}, and every run stores its fully resolved
configuration. \todo{Add the commit hash, hardware, total compute, and a
Zenodo DOI for the run archive.}

\section*{Acknowledgements}
\todo{Acknowledgements.}

\bibliographystyle{plainnat}
\bibliography{refs}

% =====================================================================
\appendix
\section{Note on an Earlier Code Release}\label{app:errata}

An earlier version of this repository (v0) reported that a frame-stack PPO
agent solved random $11\times11$ mazes, with an average return of about 9.6.
An audit found the following:
(i) the maze generator never reset its grid, so each environment instance
kept a single layout for its whole lifetime. Training with 8 parallel
environments therefore used 8 fixed layouts, and evaluation used a ninth;
(ii) the visualisation script drew a different maze from the one the
policy acted in;
(iii) v0's own evaluation log on its unseen layout shows success rates of
5--20\% and mean returns near $-90$;
(iv) the distance-shaping term was not potential-based.
All of these are fixed in the version described here, and v0's claims are
withdrawn. The \emph{legacy}/8-layout condition in \cref{fig:reward} reproduces v0's
training setup under the corrected evaluation protocol.

\section{Additional Results}
\todo{Deterministic-policy versions of the figures, wall-hit rates, and
example trajectories (\texttt{scripts/render.py}).}

\end{document}
